\section{Comparison with mathlib homology}

The definition \code{toMathlibChainComplex} sends an
\code{IntChainComplex} to
\[
  \code{ChainComplex}\;(\code{ModuleCat}\ \Z)\;\mathbb{N}.
\]
Its object in degree $n$ is the finite function module
$\Z^{\rank C_n}$, and its differential from degree $n+1$ to $n$ is the
linear map induced by \code{boundary n}.  The homological-complex chain law
is a direct consequence of \code{boundary\_squared}; the theorem
\code{toMathlibChainComplex\_d} exposes the differential equation.

The intrinsic quotient is compared to mathlib homology through short
complexes.  In positive degree the incoming, middle, and outgoing objects
agree up to the index equalities used by mathlib's complex shape.  Degree
zero needs a genuine comparison: the project's outgoing map has codomain
$0$, while mathlib's canonical short complex uses the zero endomorphism
$C_0\to C_0$.  The morphism \code{degreeShortComplexToSc} has identity
middle component and the unique map $0\to C_0$ as its third component.
Its first component is epic, its middle component is an isomorphism, and its
third component is monic, including the empty-domain proof in degree zero.
Mathlib's homology-map theorem then yields \code{degreeHomologyIso}.

Composing this categorical isomorphism with the definitional quotient
identification gives
\[
  \code{mathlibHomologyEquiv}\ n :
  H_n \simeq_{\Z}
  \code{toMathlibChainComplex.homology}\ n.
\]
The definition \code{mathlibHomologyDecomposition} transports the complete
cyclic-plus-free decomposition across this equivalence.

\subsection{Normalized Moore realization}

The finite normalized data of Section~\ref{sec:normalized-input} records
nondegenerate face outcomes and the equation $D^2=0$.  These facts do not
determine degeneracy maps or all simplicial identities, so they are
insufficient to construct an arbitrary simplicial set.  The structure
\code{NormalizedMooreComparison} makes the missing boundary explicit.  It
contains a genuine mathlib simplicial set, bijections from the finite indices
to its nondegenerate simplices, degreewise basis equivalences, agreement of
\code{some} faces, proof that \code{none} faces are degenerate, agreement on
canonical generators, and differential commutation.

Given this witness, \code{chainIso} identifies the explicit matrix complex
with mathlib's nondegenerate normalized chain complex.  The canonical
splitting of the free module on a simplicial set then connects that complex
to the normalized Moore complex in the Karoubi envelope.  Both endpoints are
in the image of the fully faithful Karoubi embedding, so
\code{normalizedMooreChainIso} reflects the result to an actual isomorphism
of chain complexes.  Applying mathlib's homology functor and the explicit
decomposition produces \code{normalizedChainHomologyDecomposition} and
\code{normalizedMooreHomologyDecomposition}.

This parameterization separates two claims cleanly: the finite matrix
homology theorem is unconditional, while transport to a particular
simplicial model requires a checked realization of that model.
